% Include from the existing main document with % Include from the existing main document with % Include from the existing main document with % Include from the existing main document with \input{roadmap}.
% This section uses core LaTeX only and inherits the parent document style.
\section{Project Execution Roadmap}

This roadmap turns the proposed study into team assignments and a reusable
Rust tool. The research question and graph-level Hodge method remain those
of the proposal. Detailed tasks, dependencies, file contracts, and acceptance
checks are recorded in \texttt{ROADMAP.md}. Environmental-response evidence
prioritizes genes; without patient exposure measurements, it does not establish
that an exposure caused a splicing change.

\subsection{Team Workstreams}

The genetics team will collect susceptibility evidence, normalize stable gene
identifiers, retain source provenance, and export a checked gene table. The
environment team will record exposure and tissue context, build the corresponding
response-gene table, and prepare risk, environmental, and overlap memberships.
These two tasks can proceed in parallel.

The data team will obtain the selected public endometriosis and control data,
verify tissue and sample metadata, obtain or derive junction counts, and export
quality-checked standard tables. The mathematics team will specify graph
orientation, incidence matrices, normalization, projections, numerical tolerances,
and synthetic expected answers. Dataset preparation and synthetic validation can
proceed simultaneously.

The Rust team will build input validation, sample and gene grouping, graph and
matrix construction, decomposition, comparison, ranking, and export modules.
Parser development can begin with synthetic inputs before real data are ready.
The statistics team will compare per-sample measurements, report effect sizes,
validate permutation tests, correct for multiple testing, and compare simpler
baselines. The presentation team will prepare pipeline and graph figures,
reviewed biological explanations, and a small live demonstration. Team members
may share roles, but each deliverable will have an owner and a reviewer.

\subsection{Development Phases}

The phase numbers identify dependencies rather than requiring every task to run
in sequence. Schemas and mathematical checks must be agreed before real-data
integration; biological interpretation follows validated statistical outputs.

\begin{center}
\small
\begin{tabular}{@{}p{0.11\linewidth}p{0.40\linewidth}p{0.39\linewidth}@{}}
\hline
Phase & Main work & Deliverable \\
\hline
1 & Organize sources, roles, and interfaces & Assigned tasks and data contract \\
2 & Build gene evidence and memberships & Risk, environment, candidate tables \\
3 & Prepare data and verify metadata & Standard junction tables and QC \\
4 & Validate synthetic mathematics & Known answers and numerical checks \\
5 & Implement Rust core & Synthetic end-to-end CLI \\
6 & Integrate real junction data & Checked real-data run \\
7 & Compare disease and control & Per-sample values and differences \\
8 & Test effects and compare baselines & Statistical and baseline outputs \\
9 & Interpret supported findings & Evidence-linked explanations \\
10 & Complete tool usability & Documented, tested Rust release \\
11 & Rehearse small live example & Fast demo and offline backup \\
12 & Assemble report and presentation & Figures, slides, reproducible methods \\
13 & Optional proposal extensions & Separately validated extensions \\
\hline
\end{tabular}
\end{center}

Core phases are required. Performance improvements are desirable when needed;
patient-level exposure analysis, full simplicial Hodge analysis, genome-wide
scaling, and other-disease applications remain optional extensions after the
core deliverables are ready.

\subsection{Minimum Viable Tool}

The minimum software deliverable is a reusable Rust command-line program. It
will read standardized junction counts, validate sample and gene identifiers,
construct directed splice graphs and the incidence matrix, normalize junction
usage, and calculate the gradient and cycle-space projections defined in the
proposal. It will retain per-sample measurements, compare eligible case and
control groups, rank genes, identify contributing junctions, and export both
readable summaries and machine-readable results with run parameters.

Synthetic validation will cover linear and branching trees, exon skipping,
back-splicing, noise, and changing junction weights. Checks will verify signal
reconstruction, cycle-space membership, orthogonality, and energy accounting.
A graph cycle-space component does not by itself establish a directed biological
cycle or a circular RNA molecule. Zero-support inputs and numerical failures
must receive explicit statuses.

The software is named Splice Girl and its command is \texttt{splice-girl}.
The following commands illustrate the intended interface; these subcommands are
not yet implemented:
\begin{verbatim}
splice-girl analyze sample.tsv
splice-girl gene GENE_ID --case cases/ --control controls/
splice-girl compare --case cases/ --control controls/
splice-girl rank --genes candidate_genes.tsv \
    --case cases/ --control controls/
\end{verbatim}

The first release may analyze a small candidate set. It must remain usable with
new conforming input files and include installation instructions, input and
output specifications, example data, numerical tests, and actionable errors.
A report or notebook alone does not satisfy the software deliverable.

\subsection{Midpoint Milestone}

By midpoint, the team should have documented gene-set criteria and provenance,
candidate tables, a selected accessible dataset with reviewed metadata and a
prepared small junction subset, and stable input/output schemas. Synthetic
expected answers and mathematical validation should be ready. The Rust parser,
graph constructor, and incidence-matrix code should work, with one synthetic
example running through normalization, decomposition, and export. Final
biological results are not required at this checkpoint. Any missing data
prerequisite must remain recorded as an outstanding task.

\subsection{Final Deliverables}

The final deliverables are a documented reusable Rust CLI, integrated real
endometriosis/control data, validated per-sample measurements, group effects,
statistical testing where the sample design permits it, baseline comparisons,
and evidence-linked interpretation. Statistical outputs will retain sample
counts, exclusions, permutation settings, and multiple-testing adjustments.
Figures and final documentation will distinguish observed findings from
illustrative examples and state the limitations of the analysis.

The live demonstration will show a tiny synthetic graph, run its decomposition,
and then run one reviewed real gene through the case/control command. The team
will explain the nodes, junctions, cycle-space measurement, group difference,
and contributing junctions. Each live command should finish in seconds during
rehearsal. Larger analyses will be precomputed and clearly labeled. Saved output,
matching figures, and a short recording will provide an offline backup.

\subsection{Risks and Fallbacks}

If the overlap gene set is small, the team can retain the risk and environmental
sets separately, as already allowed in the proposal. If junction tables are
unavailable, the data team can derive them from available sequencing files with
established software. Poor coverage or unsuitable metadata must be documented;
insufficient biological replication supports descriptive comparisons, with
inference explicitly unavailable.

If the full dataset is too large, the team will use the proposal's small-gene-set
MVP while preserving the graph-level method and real-data comparison. If the
Hodge summary offers little advantage over simpler measurements, that outcome
will be reported honestly. Full simplicial analysis will remain deferred unless
biologically meaningful higher-order structures can be justified. A synthetic-only
demonstration is useful progress but does not complete the real-data milestone.

Exact dataset accessions, exposure scope, graph representation, support filters,
normalization settings, numerical library and tolerances, statistical settings,
and final CLI details will be recorded as implementation decisions. They do not
reopen the approved research question or replace its mathematical method.
.
% This section uses core LaTeX only and inherits the parent document style.
\section{Project Execution Roadmap}

This roadmap turns the proposed study into team assignments and a reusable
Rust tool. The research question and graph-level Hodge method remain those
of the proposal. Detailed tasks, dependencies, file contracts, and acceptance
checks are recorded in \texttt{ROADMAP.md}. Environmental-response evidence
prioritizes genes; without patient exposure measurements, it does not establish
that an exposure caused a splicing change.

\subsection{Team Workstreams}

The genetics team will collect susceptibility evidence, normalize stable gene
identifiers, retain source provenance, and export a checked gene table. The
environment team will record exposure and tissue context, build the corresponding
response-gene table, and prepare risk, environmental, and overlap memberships.
These two tasks can proceed in parallel.

The data team will obtain the selected public endometriosis and control data,
verify tissue and sample metadata, obtain or derive junction counts, and export
quality-checked standard tables. The mathematics team will specify graph
orientation, incidence matrices, normalization, projections, numerical tolerances,
and synthetic expected answers. Dataset preparation and synthetic validation can
proceed simultaneously.

The Rust team will build input validation, sample and gene grouping, graph and
matrix construction, decomposition, comparison, ranking, and export modules.
Parser development can begin with synthetic inputs before real data are ready.
The statistics team will compare per-sample measurements, report effect sizes,
validate permutation tests, correct for multiple testing, and compare simpler
baselines. The presentation team will prepare pipeline and graph figures,
reviewed biological explanations, and a small live demonstration. Team members
may share roles, but each deliverable will have an owner and a reviewer.

\subsection{Development Phases}

The phase numbers identify dependencies rather than requiring every task to run
in sequence. Schemas and mathematical checks must be agreed before real-data
integration; biological interpretation follows validated statistical outputs.

\begin{center}
\small
\begin{tabular}{@{}p{0.11\linewidth}p{0.40\linewidth}p{0.39\linewidth}@{}}
\hline
Phase & Main work & Deliverable \\
\hline
1 & Organize sources, roles, and interfaces & Assigned tasks and data contract \\
2 & Build gene evidence and memberships & Risk, environment, candidate tables \\
3 & Prepare data and verify metadata & Standard junction tables and QC \\
4 & Validate synthetic mathematics & Known answers and numerical checks \\
5 & Implement Rust core & Synthetic end-to-end CLI \\
6 & Integrate real junction data & Checked real-data run \\
7 & Compare disease and control & Per-sample values and differences \\
8 & Test effects and compare baselines & Statistical and baseline outputs \\
9 & Interpret supported findings & Evidence-linked explanations \\
10 & Complete tool usability & Documented, tested Rust release \\
11 & Rehearse small live example & Fast demo and offline backup \\
12 & Assemble report and presentation & Figures, slides, reproducible methods \\
13 & Optional proposal extensions & Separately validated extensions \\
\hline
\end{tabular}
\end{center}

Core phases are required. Performance improvements are desirable when needed;
patient-level exposure analysis, full simplicial Hodge analysis, genome-wide
scaling, and other-disease applications remain optional extensions after the
core deliverables are ready.

\subsection{Minimum Viable Tool}

The minimum software deliverable is a reusable Rust command-line program. It
will read standardized junction counts, validate sample and gene identifiers,
construct directed splice graphs and the incidence matrix, normalize junction
usage, and calculate the gradient and cycle-space projections defined in the
proposal. It will retain per-sample measurements, compare eligible case and
control groups, rank genes, identify contributing junctions, and export both
readable summaries and machine-readable results with run parameters.

Synthetic validation will cover linear and branching trees, exon skipping,
back-splicing, noise, and changing junction weights. Checks will verify signal
reconstruction, cycle-space membership, orthogonality, and energy accounting.
A graph cycle-space component does not by itself establish a directed biological
cycle or a circular RNA molecule. Zero-support inputs and numerical failures
must receive explicit statuses.

The software is named Splice Girl and its command is \texttt{splice-girl}.
The following commands illustrate the intended interface; these subcommands are
not yet implemented:
\begin{verbatim}
splice-girl analyze sample.tsv
splice-girl gene GENE_ID --case cases/ --control controls/
splice-girl compare --case cases/ --control controls/
splice-girl rank --genes candidate_genes.tsv \
    --case cases/ --control controls/
\end{verbatim}

The first release may analyze a small candidate set. It must remain usable with
new conforming input files and include installation instructions, input and
output specifications, example data, numerical tests, and actionable errors.
A report or notebook alone does not satisfy the software deliverable.

\subsection{Midpoint Milestone}

By midpoint, the team should have documented gene-set criteria and provenance,
candidate tables, a selected accessible dataset with reviewed metadata and a
prepared small junction subset, and stable input/output schemas. Synthetic
expected answers and mathematical validation should be ready. The Rust parser,
graph constructor, and incidence-matrix code should work, with one synthetic
example running through normalization, decomposition, and export. Final
biological results are not required at this checkpoint. Any missing data
prerequisite must remain recorded as an outstanding task.

\subsection{Final Deliverables}

The final deliverables are a documented reusable Rust CLI, integrated real
endometriosis/control data, validated per-sample measurements, group effects,
statistical testing where the sample design permits it, baseline comparisons,
and evidence-linked interpretation. Statistical outputs will retain sample
counts, exclusions, permutation settings, and multiple-testing adjustments.
Figures and final documentation will distinguish observed findings from
illustrative examples and state the limitations of the analysis.

The live demonstration will show a tiny synthetic graph, run its decomposition,
and then run one reviewed real gene through the case/control command. The team
will explain the nodes, junctions, cycle-space measurement, group difference,
and contributing junctions. Each live command should finish in seconds during
rehearsal. Larger analyses will be precomputed and clearly labeled. Saved output,
matching figures, and a short recording will provide an offline backup.

\subsection{Risks and Fallbacks}

If the overlap gene set is small, the team can retain the risk and environmental
sets separately, as already allowed in the proposal. If junction tables are
unavailable, the data team can derive them from available sequencing files with
established software. Poor coverage or unsuitable metadata must be documented;
insufficient biological replication supports descriptive comparisons, with
inference explicitly unavailable.

If the full dataset is too large, the team will use the proposal's small-gene-set
MVP while preserving the graph-level method and real-data comparison. If the
Hodge summary offers little advantage over simpler measurements, that outcome
will be reported honestly. Full simplicial analysis will remain deferred unless
biologically meaningful higher-order structures can be justified. A synthetic-only
demonstration is useful progress but does not complete the real-data milestone.

Exact dataset accessions, exposure scope, graph representation, support filters,
normalization settings, numerical library and tolerances, statistical settings,
and final CLI details will be recorded as implementation decisions. They do not
reopen the approved research question or replace its mathematical method.
.
% This section uses core LaTeX only and inherits the parent document style.
\section{Project Execution Roadmap}

This roadmap turns the proposed study into team assignments and a reusable
Rust tool. The research question and graph-level Hodge method remain those
of the proposal. Detailed tasks, dependencies, file contracts, and acceptance
checks are recorded in \texttt{ROADMAP.md}. Environmental-response evidence
prioritizes genes; without patient exposure measurements, it does not establish
that an exposure caused a splicing change.

\subsection{Team Workstreams}

The genetics team will collect susceptibility evidence, normalize stable gene
identifiers, retain source provenance, and export a checked gene table. The
environment team will record exposure and tissue context, build the corresponding
response-gene table, and prepare risk, environmental, and overlap memberships.
These two tasks can proceed in parallel.

The data team will obtain the selected public endometriosis and control data,
verify tissue and sample metadata, obtain or derive junction counts, and export
quality-checked standard tables. The mathematics team will specify graph
orientation, incidence matrices, normalization, projections, numerical tolerances,
and synthetic expected answers. Dataset preparation and synthetic validation can
proceed simultaneously.

The Rust team will build input validation, sample and gene grouping, graph and
matrix construction, decomposition, comparison, ranking, and export modules.
Parser development can begin with synthetic inputs before real data are ready.
The statistics team will compare per-sample measurements, report effect sizes,
validate permutation tests, correct for multiple testing, and compare simpler
baselines. The presentation team will prepare pipeline and graph figures,
reviewed biological explanations, and a small live demonstration. Team members
may share roles, but each deliverable will have an owner and a reviewer.

\subsection{Development Phases}

The phase numbers identify dependencies rather than requiring every task to run
in sequence. Schemas and mathematical checks must be agreed before real-data
integration; biological interpretation follows validated statistical outputs.

\begin{center}
\small
\begin{tabular}{@{}p{0.11\linewidth}p{0.40\linewidth}p{0.39\linewidth}@{}}
\hline
Phase & Main work & Deliverable \\
\hline
1 & Organize sources, roles, and interfaces & Assigned tasks and data contract \\
2 & Build gene evidence and memberships & Risk, environment, candidate tables \\
3 & Prepare data and verify metadata & Standard junction tables and QC \\
4 & Validate synthetic mathematics & Known answers and numerical checks \\
5 & Implement Rust core & Synthetic end-to-end CLI \\
6 & Integrate real junction data & Checked real-data run \\
7 & Compare disease and control & Per-sample values and differences \\
8 & Test effects and compare baselines & Statistical and baseline outputs \\
9 & Interpret supported findings & Evidence-linked explanations \\
10 & Complete tool usability & Documented, tested Rust release \\
11 & Rehearse small live example & Fast demo and offline backup \\
12 & Assemble report and presentation & Figures, slides, reproducible methods \\
13 & Optional proposal extensions & Separately validated extensions \\
\hline
\end{tabular}
\end{center}

Core phases are required. Performance improvements are desirable when needed;
patient-level exposure analysis, full simplicial Hodge analysis, genome-wide
scaling, and other-disease applications remain optional extensions after the
core deliverables are ready.

\subsection{Minimum Viable Tool}

The minimum software deliverable is a reusable Rust command-line program. It
will read standardized junction counts, validate sample and gene identifiers,
construct directed splice graphs and the incidence matrix, normalize junction
usage, and calculate the gradient and cycle-space projections defined in the
proposal. It will retain per-sample measurements, compare eligible case and
control groups, rank genes, identify contributing junctions, and export both
readable summaries and machine-readable results with run parameters.

Synthetic validation will cover linear and branching trees, exon skipping,
back-splicing, noise, and changing junction weights. Checks will verify signal
reconstruction, cycle-space membership, orthogonality, and energy accounting.
A graph cycle-space component does not by itself establish a directed biological
cycle or a circular RNA molecule. Zero-support inputs and numerical failures
must receive explicit statuses.

The software is named Splice Girl and its command is \texttt{splice-girl}.
The following commands illustrate the intended interface; these subcommands are
not yet implemented:
\begin{verbatim}
splice-girl analyze sample.tsv
splice-girl gene GENE_ID --case cases/ --control controls/
splice-girl compare --case cases/ --control controls/
splice-girl rank --genes candidate_genes.tsv \
    --case cases/ --control controls/
\end{verbatim}

The first release may analyze a small candidate set. It must remain usable with
new conforming input files and include installation instructions, input and
output specifications, example data, numerical tests, and actionable errors.
A report or notebook alone does not satisfy the software deliverable.

\subsection{Midpoint Milestone}

By midpoint, the team should have documented gene-set criteria and provenance,
candidate tables, a selected accessible dataset with reviewed metadata and a
prepared small junction subset, and stable input/output schemas. Synthetic
expected answers and mathematical validation should be ready. The Rust parser,
graph constructor, and incidence-matrix code should work, with one synthetic
example running through normalization, decomposition, and export. Final
biological results are not required at this checkpoint. Any missing data
prerequisite must remain recorded as an outstanding task.

\subsection{Final Deliverables}

The final deliverables are a documented reusable Rust CLI, integrated real
endometriosis/control data, validated per-sample measurements, group effects,
statistical testing where the sample design permits it, baseline comparisons,
and evidence-linked interpretation. Statistical outputs will retain sample
counts, exclusions, permutation settings, and multiple-testing adjustments.
Figures and final documentation will distinguish observed findings from
illustrative examples and state the limitations of the analysis.

The live demonstration will show a tiny synthetic graph, run its decomposition,
and then run one reviewed real gene through the case/control command. The team
will explain the nodes, junctions, cycle-space measurement, group difference,
and contributing junctions. Each live command should finish in seconds during
rehearsal. Larger analyses will be precomputed and clearly labeled. Saved output,
matching figures, and a short recording will provide an offline backup.

\subsection{Risks and Fallbacks}

If the overlap gene set is small, the team can retain the risk and environmental
sets separately, as already allowed in the proposal. If junction tables are
unavailable, the data team can derive them from available sequencing files with
established software. Poor coverage or unsuitable metadata must be documented;
insufficient biological replication supports descriptive comparisons, with
inference explicitly unavailable.

If the full dataset is too large, the team will use the proposal's small-gene-set
MVP while preserving the graph-level method and real-data comparison. If the
Hodge summary offers little advantage over simpler measurements, that outcome
will be reported honestly. Full simplicial analysis will remain deferred unless
biologically meaningful higher-order structures can be justified. A synthetic-only
demonstration is useful progress but does not complete the real-data milestone.

Exact dataset accessions, exposure scope, graph representation, support filters,
normalization settings, numerical library and tolerances, statistical settings,
and final CLI details will be recorded as implementation decisions. They do not
reopen the approved research question or replace its mathematical method.
.
% This section uses core LaTeX only and inherits the parent document style.
\section{Project Execution Roadmap}

This roadmap turns the proposed study into team assignments and a reusable
Rust tool. The research question and graph-level Hodge method remain those
of the proposal. Detailed tasks, dependencies, file contracts, and acceptance
checks are recorded in \texttt{ROADMAP.md}. Environmental-response evidence
prioritizes genes; without patient exposure measurements, it does not establish
that an exposure caused a splicing change.

\subsection{Team Workstreams}

The genetics team will collect susceptibility evidence, normalize stable gene
identifiers, retain source provenance, and export a checked gene table. The
environment team will record exposure and tissue context, build the corresponding
response-gene table, and prepare risk, environmental, and overlap memberships.
These two tasks can proceed in parallel.

The data team will obtain the selected public endometriosis and control data,
verify tissue and sample metadata, obtain or derive junction counts, and export
quality-checked standard tables. The mathematics team will specify graph
orientation, incidence matrices, normalization, projections, numerical tolerances,
and synthetic expected answers. Dataset preparation and synthetic validation can
proceed simultaneously.

The Rust team will build input validation, sample and gene grouping, graph and
matrix construction, decomposition, comparison, ranking, and export modules.
Parser development can begin with synthetic inputs before real data are ready.
The statistics team will compare per-sample measurements, report effect sizes,
validate permutation tests, correct for multiple testing, and compare simpler
baselines. The presentation team will prepare pipeline and graph figures,
reviewed biological explanations, and a small live demonstration. Team members
may share roles, but each deliverable will have an owner and a reviewer.

\subsection{Development Phases}

The phase numbers identify dependencies rather than requiring every task to run
in sequence. Schemas and mathematical checks must be agreed before real-data
integration; biological interpretation follows validated statistical outputs.

\begin{center}
\small
\begin{tabular}{@{}p{0.11\linewidth}p{0.40\linewidth}p{0.39\linewidth}@{}}
\hline
Phase & Main work & Deliverable \\
\hline
1 & Organize sources, roles, and interfaces & Assigned tasks and data contract \\
2 & Build gene evidence and memberships & Risk, environment, candidate tables \\
3 & Prepare data and verify metadata & Standard junction tables and QC \\
4 & Validate synthetic mathematics & Known answers and numerical checks \\
5 & Implement Rust core & Synthetic end-to-end CLI \\
6 & Integrate real junction data & Checked real-data run \\
7 & Compare disease and control & Per-sample values and differences \\
8 & Test effects and compare baselines & Statistical and baseline outputs \\
9 & Interpret supported findings & Evidence-linked explanations \\
10 & Complete tool usability & Documented, tested Rust release \\
11 & Rehearse small live example & Fast demo and offline backup \\
12 & Assemble report and presentation & Figures, slides, reproducible methods \\
13 & Optional proposal extensions & Separately validated extensions \\
\hline
\end{tabular}
\end{center}

Core phases are required. Performance improvements are desirable when needed;
patient-level exposure analysis, full simplicial Hodge analysis, genome-wide
scaling, and other-disease applications remain optional extensions after the
core deliverables are ready.

\subsection{Minimum Viable Tool}

The minimum software deliverable is a reusable Rust command-line program. It
will read standardized junction counts, validate sample and gene identifiers,
construct directed splice graphs and the incidence matrix, normalize junction
usage, and calculate the gradient and cycle-space projections defined in the
proposal. It will retain per-sample measurements, compare eligible case and
control groups, rank genes, identify contributing junctions, and export both
readable summaries and machine-readable results with run parameters.

Synthetic validation will cover linear and branching trees, exon skipping,
back-splicing, noise, and changing junction weights. Checks will verify signal
reconstruction, cycle-space membership, orthogonality, and energy accounting.
A graph cycle-space component does not by itself establish a directed biological
cycle or a circular RNA molecule. Zero-support inputs and numerical failures
must receive explicit statuses.

The software is named Splice Girl and its command is \texttt{splice-girl}.
The following commands illustrate the intended interface; these subcommands are
not yet implemented:
\begin{verbatim}
splice-girl analyze sample.tsv
splice-girl gene GENE_ID --case cases/ --control controls/
splice-girl compare --case cases/ --control controls/
splice-girl rank --genes candidate_genes.tsv \
    --case cases/ --control controls/
\end{verbatim}

The first release may analyze a small candidate set. It must remain usable with
new conforming input files and include installation instructions, input and
output specifications, example data, numerical tests, and actionable errors.
A report or notebook alone does not satisfy the software deliverable.

\subsection{Midpoint Milestone}

By midpoint, the team should have documented gene-set criteria and provenance,
candidate tables, a selected accessible dataset with reviewed metadata and a
prepared small junction subset, and stable input/output schemas. Synthetic
expected answers and mathematical validation should be ready. The Rust parser,
graph constructor, and incidence-matrix code should work, with one synthetic
example running through normalization, decomposition, and export. Final
biological results are not required at this checkpoint. Any missing data
prerequisite must remain recorded as an outstanding task.

\subsection{Final Deliverables}

The final deliverables are a documented reusable Rust CLI, integrated real
endometriosis/control data, validated per-sample measurements, group effects,
statistical testing where the sample design permits it, baseline comparisons,
and evidence-linked interpretation. Statistical outputs will retain sample
counts, exclusions, permutation settings, and multiple-testing adjustments.
Figures and final documentation will distinguish observed findings from
illustrative examples and state the limitations of the analysis.

The live demonstration will show a tiny synthetic graph, run its decomposition,
and then run one reviewed real gene through the case/control command. The team
will explain the nodes, junctions, cycle-space measurement, group difference,
and contributing junctions. Each live command should finish in seconds during
rehearsal. Larger analyses will be precomputed and clearly labeled. Saved output,
matching figures, and a short recording will provide an offline backup.

\subsection{Risks and Fallbacks}

If the overlap gene set is small, the team can retain the risk and environmental
sets separately, as already allowed in the proposal. If junction tables are
unavailable, the data team can derive them from available sequencing files with
established software. Poor coverage or unsuitable metadata must be documented;
insufficient biological replication supports descriptive comparisons, with
inference explicitly unavailable.

If the full dataset is too large, the team will use the proposal's small-gene-set
MVP while preserving the graph-level method and real-data comparison. If the
Hodge summary offers little advantage over simpler measurements, that outcome
will be reported honestly. Full simplicial analysis will remain deferred unless
biologically meaningful higher-order structures can be justified. A synthetic-only
demonstration is useful progress but does not complete the real-data milestone.

Exact dataset accessions, exposure scope, graph representation, support filters,
normalization settings, numerical library and tolerances, statistical settings,
and final CLI details will be recorded as implementation decisions. They do not
reopen the approved research question or replace its mathematical method.
